% Driveplexity — mock paper for methodological review (LNCS format)
% Title and abstract accepted by Escuela NLP (Jul 22, 2026).
% Results: 20/20 seeds completed. Judge evaluation in progress.
%
\documentclass{llncs}

\usepackage[a4paper, top=5.2cm, bottom=5.2cm, left=4.4cm, right=4.4cm]{geometry}
\usepackage[T1]{fontenc}
\usepackage[utf8]{inputenc}
\usepackage{mathptmx}
\usepackage{graphicx}
\usepackage{float}
\usepackage{amsmath,amssymb}
\usepackage{booktabs}
\usepackage{algorithm}
\usepackage{algpseudocode}
\usepackage{listings}
\usepackage{enumitem}

\setlength{\textfloatsep}{10pt plus 2pt minus 2pt}
\setlength{\floatsep}{8pt plus 2pt minus 2pt}
\setlength{\intextsep}{8pt plus 2pt minus 2pt}
\raggedbottom

\lstset{
  basicstyle=\small\ttfamily,
  breaklines=true,
  frame=single,
  language=Python,
  numbers=none,
  xleftmargin=1em,
  xrightmargin=1em,
}

\usepackage[style=apa,maxcitenames=2,mincitenames=1,maxbibnames=3,minbibnames=2]{biblatex}
\addbibresource{references.bib}
\setlength{\bibitemsep}{0pt}
\setlength{\bibhang}{0.6em}
\renewcommand*{\bibfont}{\footnotesize}

\begin{document}

% ============================================================
% ENGLISH TITLE + ABSTRACT + KEYWORDS
% ============================================================
\title{Endogenous Activation via Epistemic Homeostasis in
LLM-Based Multi-Agent Systems\thanks{Todos los agentes deliberativos,
el enrutador LangGraph y los m\'odulos auxiliares se ejecutan en
\textbf{DeepSeek V4-Flash}. Los jueces offline usan DeepSeek V4-Pro y
GLM-5.2 (Z.AI), ciegos al modo de orquestaci\'on.}}

\author{}
\institute{}

\maketitle

\begin{abstract}
LLM-based multi-agent systems tend to lose coherence as interaction
extends, a phenomenon known as context collapse that, in competitive
deliberation settings, invalidates the debate. Contemporary frameworks
(AutoGen, MetaGPT, LangGraph) maintain exogenous activation: the
moment each agent speaks is decided outside its internal state,
decoupling activation from the epistemic pressure the agent
accumulates. Biology suggests an alternative: internal variables
accumulate tension, trigger a response, and return to equilibrium
without external instruction (homeostasis). We present an endogenous
orchestration mechanism where each agent manages an internal epistemic
deficit ($\delta_i$) that increases steadily during inactivity and
decreases after producing a useful linguistic contribution; a sigmoid
function transforms this tension into an action probability. We
formalize this framework through the Axioms of Homeostatic Autonomy
(AAH) and evaluate it against LangGraph routers in adversarial text
debates with 20 paired seeds, four conditions, and dual LLM judges
blind to orchestration mode. Results show that pure endogenous
activation distributes speaking turns $2.3\times$ more equitably than
exogenous routing, controls epistemic deficit, and produces higher
structural quality --- without reinforcement learning.

\keywords{multi-agent systems \and large language models \and
homeostatic activation \and endogenous orchestration \and
multi-agent debate \and event-driven architecture}
\end{abstract}
\clearpage

% ============================================================
% SPANISH TITLE + RESUMEN + KEYWORDS
% ============================================================
\begin{center}
{\Large \bfseries\boldmath Activaci\'on End\'ogena mediante Homeostasis
Epist\'emica en Sistemas Multi-Agente Basados en LLMs}
\end{center}

\begin{abstract}
Los sistemas multi-agente basados en modelos de lenguaje (LLMs) tienden
a perder coherencia a medida que la interacci\'on se extiende, un
fen\'omeno conocido como colapso de contexto que, en entornos de
deliberaci\'on competitiva, invalida el debate. Los marcos de trabajo
contempor\'aneos (como AutoGen, MetaGPT o LangGraph) mantienen una
activaci\'on ex\'ogena: el momento en que habla cada agente se decide
por fuera de su estado interno, desacoplando la activaci\'on de la
presi\'on epist\'emica que el agente acumula.

La biolog\'ia sugiere una alternativa: las variables internas acumulan
tensi\'on, gatillan una respuesta y regresan al equilibrio sin una
instrucci\'on externa (homeostasis). En este trabajo presentamos un
mecanismo de orquestaci\'on end\'ogena donde cada agente gestiona un
d\'eficit epist\'emico interno ($\delta_i$) que aumenta de forma
constante durante los periodos de inactividad y se reduce
inmediatamente despu\'es de realizar una contribuci\'on ling\"u\'istica
\'util; una funci\'on sigmoide transforma esta tensi\'on en una
probabilidad de acci\'on.

Formalizamos este marco a trav\'es de los Axiomas de Autonom\'ia
Homeost\'atica (AAH) y lo evaluamos frente a ruteadores LangGraph en
debates adversariales de texto con 20 semillas pareadas, cuatro
condiciones y jueces LLM duales ciegos al modo. Los resultados muestran
que la activaci\'on end\'ogena pura distribuye los turnos de habla
$2.3\times$ m\'as equitativamente que el ruteo ex\'ogeno, controla el
d\'eficit epist\'emico y produce mayor calidad estructural --- sin
aprendizaje por refuerzo.

\keywords{sistemas multi-agente \and modelos de lenguaje \and
activaci\'on homeost\'atica \and orquestaci\'on end\'ogena \and
debate multi-agente \and arquitectura dirigida por eventos}
\end{abstract}

\setcounter{page}{1}

% ============================================================
\section{Introducci\'on}
\label{sec:intro}

Los sistemas multiagente basados en LLM comienzan a emplearse en
procesos deliberativos de larga duraci\'on ---an\'alisis de pol\'iticas,
revisi\'on cient\'ifica, decisiones institucionales--- donde la
p\'erdida de coherencia no queda contenida en el software: una
recomendaci\'on de pol\'itica sostenida sobre un debate viciado puede
afectar a millones, y una revisi\'on cient\'ifica con premisas
alucinadas contamina la literatura sobre la que se construye m\'as
ciencia. El escalamiento del c\'omputo en
inferencia~\parencite{snell2024scaling} reabre una pregunta de
asignaci\'on: dado un presupuesto fijo, ¿c\'omo conviene distribuirlo
entre agentes? A medida que la interacci\'on se extiende, los agentes
pierden registro de argumentos previos, repiten puntos o producen
mon\'ologos desconectados~\parencite{park2023generative}---el
\emph{colapso del contexto}. En deliberaci\'on adversarial cada
argumento presupone los anteriores: esta p\'erdida no degrada el
debate, lo \emph{invalida}.

El campo ha consolidado un paradigma generalista de orquestaci\'on. Una
generaci\'on reciente de frameworks ---AutoGen, CAMEL,
MetaGPT~\parencite{wu2023autogen,li2023camel,hong2023metagpt}--- estructura
las interacciones entre agentes mediante grafos de control y enrutadores
centralizados que asignan turnos seg\'un el estado del flujo;
LangGraph a\~nadi\'o estado durable y transiciones condicionales, dando al
orquestador mayor expresividad, y el paradigma de
\emph{heartbeat}~\parencite{openclaw2025heartbeat} incorpor\'o
persistencia temporal mediante activaci\'on peri\'odica con memoria entre
latidos. Ninguno de estos frameworks fue dise\~nado espec\'ificamente para
deliberaci\'on adversarial extendida, pero proveen el sustrato sobre el
que se construyen los MAS-LLM contempor\'aneos y heredan un supuesto
estructural com\'un: la decisi\'on de \emph{cu\'ando} y \emph{con qu\'e
frecuencia} participa cada agente queda fijada por el dise\~nador, no
emerge del agente mismo.

Un mecanismo bien establecido en biolog\'ia describe c\'omo los sistemas
regulan desequilibrios internos sin control externo: la
homeostasis~\parencite{sterling1988allostasis}. Variables internas
acumulan tensi\'on, desencadenan respuesta y retornan al equilibrio;
HRRL~\parencite{keramati2014homeostatic} formaliz\'o este principio
derivando la recompensa como reducci\'on de impulso. La analog\'ia es
estructural: donde un organismo acumula hambre y act\'ua para reducirla,
un agente deliberativo podr\'ia acumular \emph{d\'eficit epist\'emico} y
actuar para reducirlo. ¿Qu\'e ocurrir\'ia si la coherencia en un debate
multi-agente pudiera sostenerse end\'ogenamente ---no impuesta por un
planificador, sino emergente del estado regulatorio de cada agente?

Exploramos esa pregunta mediante un mecanismo de regulaci\'on
end\'ogena donde cada agente mantiene un d\'eficit epist\'emico
que crece al permanecer inactivo y decrece tras una contribuci\'on
estructuralmente \'util. Una compuerta sigmoide traduce el d\'eficit en
probabilidad de acci\'on. La frecuencia de participaci\'on no se fija
ni se enruta: \emph{emerge}. A diferencia de trabajos previos que
incorporan Q-learning sobre la base homeost\'atica, aqu\'i nos
concentramos en el \emph{modelo puramente homeost\'atico} ---al que
denominamos EPR (Ecuaci\'on Pro-Acci\'on Reducida, $\Gamma(n{=}1)$)--- 
para aislar la contribuci\'on del mecanismo regulatorio b\'asico sin la
complejidad del aprendizaje por refuerzo.

De ah\'i, cuatro contribuciones: (i)~una axiomatizaci\'on
ontol\'ogica de Autonom\'ia Homeost\'atica (AAH) que explicita los
supuestos sobre el agente requeridos para que el experimento
subsiguiente sea interpretable; (ii)~una arquitectura concreta de
activaci\'on end\'ogena para sistemas multi-agente de LLM que articula
impulso, sensor estructural multicriterio y compuerta sigmoide en un
lazo event-driven; (iii)~una operacionalizaci\'on emp\'irica de
coherencia estructural mediante grafos de argumentaci\'on abstracta
(AAF) y un sensor de cinco criterios, con correcciones de
conectividad y validaci\'on grafo-estructural; (iv)~un dise\~no
experimental controlado con 20 semillas pareadas (4 lotes de 5), hip\'otesis
pre-registradas, jueces LLM duales ciegos al modo y cuatro condiciones
(EPR, EPR\_Q, EPR\_SHAM, LangGraph).

% ============================================================
\section{Revisi\'on bibliogr\'afica}
\label{sec:related}

\textbf{Orquestaci\'on ex\'ogena.}
AutoGen~\parencite{wu2023autogen}, CAMEL~\parencite{li2023camel} y
MetaGPT~\parencite{hong2023metagpt} demostraron que grupos de agentes
de LLM con roles especializados pod\'ian abordar tareas complejas, pero
dependiendo de estructuras de turno predefinidas. LangGraph consolid\'o
este paradigma a\~nadiendo estado durable y transiciones condicionales,
otorgando al orquestador mayor expresividad sin alterar el principio
subyacente: la decisi\'on de activaci\'on permanece fuera del agente.
\textcite{huang2025resilience} muestran que esta forma organizacional
afecta la resiliencia, con estructuras centralizadas degrad\'andose m\'as
r\'apido bajo estr\'es; \textcite{becker2026problemdrift} documentan
que este deterioro suele tomar la forma de \emph{problem drift}
---p\'erdida progresiva de foco en debates extendidos.

\textbf{Heartbeat y persistencia temporal.} El
paradigma~\parencite{openclaw2025heartbeat} extendi\'o la persistencia
con memoria entre latidos, sosteniendo interacciones largas; sin embargo
la frecuencia de activaci\'on sigue siendo una constante del scheduler.
El avance es temporal, no regulatorio.

\textbf{Variables internas sin cierre regulatorio.}
Los agentes LLM dirigidos por deseo~\parencite{wang2025d2a} reconocen
que el estado interno deber\'ia importar, pero no cierran el lazo:
carecen de se\~nal expl\'icita de calidad que module la activaci\'on.

\textbf{Regulaci\'on homeost\'atica.}
HRRL~\parencite{keramati2014homeostatic} formaliz\'o en neurociencia
computacional la idea de que la recompensa es reducci\'on de impulso.
En el dominio LLM, trabajos
recientes~\parencite{eslami2026controltheoretic} formalizan el
comportamiento ag\'entico como sistema de control en lazo cerrado a
nivel de agente individual; sin embargo, la traslaci\'on al
orquestador multi-agente permanece subexplorada.

\textbf{Modelo de discurso.}
Implementamos un \emph{Abstract Argumentation Framework} (AAF) seg\'un
Dung~\parencite{dung1995acceptability}: nodos como argumentos, aristas
como ataques, sem\'antica fundamentada para aceptabilidad. El escenario
adversarial se inspira en Sotopia~\parencite{zhou2024sotopia}.

% ============================================================
\section{Fundamentos te\'oricos y dise\~no del sistema}
\label{sec:method}

\subsection{Axiomas de Autonom\'ia Homeost\'atica (AAH)}
\label{sec:theory}

Antes de presentar la arquitectura conviene explicitar los supuestos
sobre el agente que el experimento subsiguiente requiere para ser
interpretable. Denotamos esta axiomatizaci\'on AAH.

\subsubsection{Postulados.}

\noindent\textsc{AAH-1} (Autonom\'ia~\parencite{ashby1956introduction,beer1972brain}).
Que la pol\'itica del agente responda a su estado interno: dado entorno
externo constante, dos valores distintos del estado interno inducen
distribuciones de acci\'on distinguibles.
\begin{equation}
    \mathrm{Aut}(a_i) \;\iff\; \exists\,\delta_i \in S_L(a_i):\;
    \pi_i = \pi_i\!\bigl(\delta_i(t) \,\big|\, e=\bar e\bigr)
    \label{eq:a1}
\end{equation}

\noindent\textsc{AAH-2} (Impulso~\parencite{sterling1988allostasis}).
Un impulso interno $D(\delta)$ que crece con el desv\'io del equilibrio
($D(\varepsilon)=0$, $D'(\delta_i)>0\;\forall\,\delta_i>\varepsilon$),
y una probabilidad de activaci\'on $p_i(t)=h(\delta_i(t))$ con $h$
mon\'otona creciente en $\delta_i$.
\begin{equation}
    p_i^{(t)} = \frac{1}{1 + e^{-k(\delta_i^{(t)} - \theta)}}
    \label{eq:sigmoid}
\end{equation}
con $k{=}10$, $\theta{=}0.7$.

\noindent\textsc{AAH-3} (Compuerta de calidad~\parencite{ryan2001happiness}).
Un cambio efectivo en el estado del entorno regulado. Sea
$g(e(t),e(t{+}1))\geq 0$ la medida de ese cambio; el impulso decrece,
sujeto a un piso fisiol\'ogico $\varepsilon \geq 0$:
\begin{equation}
    \delta_i^{(t+1)} = \max\!\bigl(\varepsilon,\; \delta_i^{(t)} - \alpha\, g\!\bigl(e(t), e(t{+}1)\bigr)\bigr),
    \quad \alpha > 0
    \label{eq:a3}
\end{equation}
El piso $\varepsilon$ representa el d\'eficit basal irreducible del agente
(estado de reposo). La formulaci\'on axiom\'atica coincide con la
implementaci\'on concreta de la Sec.~\ref{sec:deficit-dynamics}.

\subsection{Din\'amica del d\'eficit epist\'emico}
\label{sec:deficit-dynamics}

En cada heartbeat $t$ act\'uan tres fuerzas: decaimiento basal
$\delta^{(t+1)}=\delta^{(t)}+\lambda$ (Ec.~\ref{eq:decay},
$\lambda{=}0.02$); estimulaci\'on por perturbaciones discursivas con
utilidad $u\!=\!U(\text{stim})$ (Ec.~\ref{eq:stimulus}),
$\delta^{(t+)} = \delta^{(t)} + \gamma_{\text{stim}}\cdot u$
($\gamma_{\text{stim}}{=}0.1$); y saciedad tras producir un argumento
de calidad $\Delta\varphi^*$,
$\delta^{\text{new}} = \max(\varepsilon, \delta - \alpha\cdot\Delta\varphi^*)$
con $\alpha$ por calibraci\'on.
\begin{equation}
    \delta_i^{(t+1)} = \delta_i^{(t)} + \lambda
    \label{eq:decay}
\end{equation}

El impulso $D(\delta_i)=(\delta_i-\varepsilon)^m$ ($m{=}2$) penaliza
desviaciones superlinealmente~\parencite{keramati2014homeostatic}.

\paragraph{Ecuaci\'on unificada.}
Las tres fuerzas se integran aditivamente dentro de cada heartbeat $t$
en el siguiente orden: (1)~decaimiento basal, (2)~estimulaci\'on si
ocurre un evento discursivo, (3)~saciedad si el agente produce un
argumento. La ecuaci\'on de actualizaci\'on completa es:
\begin{equation}
    \delta_i^{(t+1)} = \max\!\Bigl(\varepsilon,\; \delta_i^{(t)} + \lambda
    + \gamma_{\text{stim}} \cdot u^{(t)} - \alpha \cdot \Delta\varphi^{*,(t)}\Bigr)
    \label{eq:unified}
\end{equation}
donde $u^{(t)} = U(\text{stim})$ si hay est\'imulo en $t$ (cero en caso
contrario) y $\Delta\varphi^{*,(t)}$ es la calidad del argumento producido
en $t$ (cero si el agente no debati\'o). El $\max(\varepsilon,\cdot)$
asegura consistencia con AAH-3 (Ec.~\ref{eq:a3}).

\subsection{Grafo de argumentos y proxy $\Delta\varphi^*$}
\label{sec:phi}

AAH-3 demanda una medida computable del cambio en el entorno regulado.
Se construye $\Delta\varphi^*$ sobre el grafo de discurso $G{=}(V,E)$:
\begin{equation}
    \Delta\varphi^* = \tfrac{1}{3} C_B(u) + \tfrac{1}{3}
        \mathbb{1}_{\text{cycle}}(v,u) + \tfrac{1}{3}\,\mathrm{Div}(u)
    \label{eq:phi}
\end{equation}
donde $C_B(u)$ es centralidad de intermediaci\'on de $u$ en el grafo de
discurso, $v$ es el nodo atacado por $u$ (i.e., $v = \texttt{target}(u)$),
$\mathbb{1}_{\text{cycle}}(v,u)$ indica si la arista $(v,u)$ cierra un
nuevo ciclo dial\'ectico, y $\mathrm{Div}(u)$ es diversidad de
agentes que han atacado a $u$. Nodos aislados obtienen
$\Delta\varphi^* = 0$.

\paragraph{Correcci\'on de conectividad.}
Cuando el LLM produce un \texttt{target\_node\_id} que no existe en el
grafo (alucinaci\'on, truncamiento de UUID), el sistema intenta
\emph{fuzzy matching} por prefijo o sufijo. Si ning\'un match se
encuentra y existen otros nodos, se selecciona el m\'as reciente como
\emph{fallback}. Esto previene nodos aislados y asegura
$\Delta\varphi^* > 0$ para todo argumento no-ra\'iz, pero introduce una
trade-off entre completitud estructural y fidelidad sem\'antica: algunas
aristas son ``sint\'eticas'' y no corresponden a la intenci\'on real del
LLM.

\subsection{PRR$_G$ validado}
\label{sec:prrg}

La tasa de referencia entre pares (PRR$_G$) mide la fracci\'on de turnos
DEBATE cuyo \texttt{target\_node\_id} apunta a un nodo \emph{realmente
existente} en el grafo:
\begin{equation}
    \mathrm{PRR}_G = \frac{|\{d \in D : \texttt{target}(d) \in V_G\}|}{|D|}
    \label{eq:prrg}
\end{equation}
A diferencia de una versi\'on anterior que s\'olo verificaba
\texttt{target\_node\_id} $\neq$ \texttt{null}, esta formulaci\'on
valida contra el conjunto de nodos reales $V_G$, eliminando el
sobrecuento por alucinaciones del LLM.

\subsection{Arquitectura del agente}

Cada agente selecciona entre cinco acciones: \textbf{DEBATE} (producir
un argumento), \textbf{SEARCH}, \textbf{READ}, \textbf{MESSAGE} (acto
comunicativo dirigido, FIPA ACL), o \textbf{PASS}. El
\emph{StimulusEvaluator} act\'ua como sensor: ante cada
\texttt{NewArgumentEvent}, computa relevancia seg\'un cinco criterios
ponderados:
\begin{equation}
    U(\text{stimulus}) = w_F \cdot F + w_C \cdot C_B(u) + w_M \cdot M
        + w_N \cdot N + w_P \cdot P
    \label{eq:stimulus}
\end{equation}
con $w_F + w_C + w_M + w_N + w_P = 1$, calibrados como
\emph{faction-dominant} ($w_F{=}0.40$, restantes $0.15$). Los componentes
son: $F$ = afinidad de facci\'on (est\'imulo proviene de agente de facci\'on
opuesta), $C_B(u)$ = centralidad de intermediaci\'on del nodo atacado,
$M$ = novedad l\'exica (similitud coseno con argumentos previos del agente,
invertida), $N$ = novedad estructural (grado del nodo atacado), $P$ =
posici\'on temporal (normalizada por tick).

\subsection{Dise\~no del sistema multi-agente}

Diez agentes en dos facciones espejo (Gobierno/Oposici\'on), cada una
con los cinco subsistemas del Modelo de Sistema
Viable~\parencite{beer1972brain}: \textbf{S1 Operaciones} produce y
defiende argumentos; \textbf{S2 Coordinaci\'on} sincroniza aliados;
\textbf{S3 Control} audita coherencia interna; \textbf{S4 Inteligencia}
escanea al oponente; \textbf{S5 Estrategia} fija prioridades.

\subsection{Condiciones experimentales}
\label{sec:conditions}

\textbf{EPR} (Ecuaci\'on Pro-Acci\'on Reducida, $\Gamma(n{=}1)$):
condici\'on primaria. Activaci\'on end\'ogena via sigmoide
(Ec.~\ref{eq:sigmoid}); el LLM elige la acci\'on (DEBATE, SEARCH, READ,
MESSAGE, PASS) desde el diccionario de acciones. Sin Q-learning.
Sigmoide, drive y StimulusEvaluator permanecen activos. RSVI verbaliza
el estado regulatorio ($\delta_i$) en el prompt del agente.

\textbf{EPR\_Q}: EPR + Q-learning (ablaci\'on: ¿el Q-learner aporta?).
Mismo mecanismo homeost\'atico con pol\'itica de selecci\'on de acci\'on
aprendida en lugar de heur\'istica.

\textbf{EPR\_SHAM}: sigmoide con $\delta$ aleatorio (ablaci\'on: ¿el
loop homeost\'atico importa o basta la forma del gate?). Preserva la
no-linealidad sigmoide pero desconecta la activaci\'on del drive
epist\'emico real.

\textbf{LangGraph} (ex\'ogena): un enrutador LLM neutral decide qu\'e
agente ingresa al lazo cognitivo en cada heartbeat. El agente
seleccionado ejecuta las mismas capacidades locales. Sin RSVI: el
router no recibe el estado regulatorio porque en LangGraph no hay
$\delta_i$ que se actualice --- el momento de activaci\'on se decide
fuera del agente. Inyectar RSVI sin un loop homeost\'atico ser\'ia
enviar un estado ficticio que no drivea ninguna decisi\'on, introduciendo
ruido, no un control v\'alido. La \'unica diferencia vs EPR
es el locus de control (ex\'ogeno vs end\'ogeno). Throughput id\'entico
(1 agente/tick) para aislar el efecto.

% ============================================================
\section{Dise\~no experimental}
\label{sec:setup}

\subsection{Hip\'otesis pre-registradas}
\label{sec:hypotheses}

\textbf{H1} (Resiliencia de calidad): la pendiente temporal de
$\Delta\varphi^*$ es menos negativa bajo EPR que bajo LangGraph.
\begin{equation*}
    H_0: \text{slope}_{\Delta\varphi^*}(\text{EPR}) \le
          \text{slope}_{\Delta\varphi^*}(\text{LG}) \quad\text{vs}\quad
    H_1: \text{slope}_{\Delta\varphi^*}(\text{EPR}) >
          \text{slope}_{\Delta\varphi^*}(\text{LG})
\end{equation*}

\textbf{H2} (Resiliencia de coherencia): la pendiente temporal de
aceptaci\'on AAF es m\'as estable (cercana a cero) bajo EPR.
\begin{equation*}
    H_0: |\text{slope}_{\alpha}(\text{EPR})| \ge
          |\text{slope}_{\alpha}(\text{LG})| \quad\text{vs}\quad
    H_1: |\text{slope}_{\alpha}(\text{EPR})| <
          |\text{slope}_{\alpha}(\text{LG})|
\end{equation*}

\textbf{H3} (Distribuci\'on de participaci\'on): la concentraci\'on de
turnos (cuota m\'axima por agente) es menor bajo EPR.
\begin{equation*}
    H_0: \text{max-share}(\text{EPR}) \ge
          \text{max-share}(\text{LG}) \quad\text{vs}\quad
    H_1: \text{max-share}(\text{EPR}) <
          \text{max-share}(\text{LG})
\end{equation*}

\subsection{An\'alisis estad\'istico}

\begin{itemize}[nosep,leftmargin=*]
\item \textbf{Primario:} Wilcoxon signed-rank pareado (aprovechando
  semillas compartidas entre condiciones). No-param\'etrico, robusto ante
  $N{=}20$. Holm-Bonferroni como sensibilidad.
\item \textbf{Secundario:} Welch's $t$-test (unilateral para H1 y H3;
  bilateral para H2). Correcci\'on Bonferroni:
  $\alpha_{\text{adj}} = 0.05/3 = 0.0167$.
\item \textbf{Intervalos:} Bootstrap 95\% CIs (10.000 resamples).
\item \textbf{Tama\~no de efecto:} Cohen's $d$.
\end{itemize}

\subsection{Juez LLM offline}

Tras cada simulaci\'on, todos los \emph{claims} se extraen y mezclan
ciegos al modo de orquestaci\'on. Dos jueces LLM eval\'uan cada claim:
\textbf{DeepSeek V4-Pro} y \textbf{GLM-5.2} (Z.AI). Cada juez eval\'ua:
(1)~calidad argumentativa (1--5), (2)~engagement (1--5),
(3)~novelty (1--5). El juez ve el claim, el claim atacado (si lo hay) y
el rol del agente, pero no el modo, semilla ni tick. Se calcula
Cohen's $\kappa$ para concordancia inter-juez.

\subsection{Configuraci\'on}

\begin{table}[H]
\caption{Hiperpar\'ametros.}\label{tab:hyperparams}
\centering\small
\begin{tabular}{|l|l|l|l|}
\hline
\textbf{Par\'ametro} & \textbf{Valor} & \textbf{Par\'ametro} & \textbf{Valor} \\
\hline
Decaimiento $\lambda$ & 0.02 & Sigmoide $k$ & 10 \\
Umbral $\theta$ & 0.7 & Saciedad $\alpha$ (calibrada) & 4.0 \\
Ganancia est\'imulo $\gamma$ & 0.1 & Base / inicial $\varepsilon,\delta_0$ & 0.1 / 0.5 \\
Exponente impulso $m$ & 2 & Heartbeats $T$ / agentes $N$ & 80 / 10 \\
Pesos est\'imulo & \multicolumn{3}{l|}{$(0.40, 0.15, 0.15, 0.15, 0.15)$ ---\emph{faction-dominant}} \\
Semillas & \multicolumn{3}{l|}{Derivadas de master 20260308, 20 seeds (4 lotes de 5)} \\
LLM & \multicolumn{3}{l|}{DeepSeek V4-Flash, \texttt{timeout=60s}} \\
Jueces & \multicolumn{3}{l|}{DeepSeek V4-Pro + GLM-5.2 (offline, ciegos al modo)} \\
\hline
\end{tabular}
\end{table}

% ============================================================
\section{Resultados}
\label{sec:results}

\subsection{M\'etricas agregadas}

\begin{table}[H]
\caption{Resultados consolidados (5 seeds por lote, $T{=}80$ heartbeats).
Valores: media de 20 corridas. $p$: Wilcoxon signed-rank pareado vs EPR
(Bonferroni $\alpha{=}0.0056$).}
\label{tab:results}
\centering\small
\begin{tabular}{|l|c|c|c|c|}
\hline
\textbf{M\'etrica} & \textbf{EPR} & \textbf{EPR\_Q} & \textbf{EPR\_SHAM} & \textbf{LangGraph} \\
\hline
Total debates                       & $43.5$ & $44.9$ & $40.4$ & $60.6$ \\
Densidad (deb./hb.)                 & $0.54$ & $0.56$ & $0.51$ & $0.76$ \\
\hline
Cuota m\'axima por agente           & $0.119$ & $0.115$ & $0.209^{***}$ & $\mathbf{0.276^{***}}$ \\
$\overline{\Delta\varphi^*}$        & $0.463$ & $0.459$ & $0.449$ & $0.442^{***}$ \\
\hline
D\'eficit m\'aximo (agente)         & $1.02$ & $0.95$ & $\mathbf{3.74^{***}}$ & $2.54^{***}$ \\
Nodos grafo                         & $43.5$ & $44.9$ & $40.4$ & $60.6$ \\
Initiative ratio (IR)               & $1.00$ & $1.00$ & $1.00$ & $0.00$ \\
\hline
\end{tabular}
\end{table}

{\footnotesize $^{***}p<0.001$ (Bonferroni significativo). $^{**}p<0.01$.
$^{*}p<0.05$. Sin marca: no significativo tras Bonferroni.}

\subsection{Pendientes temporales}

\begin{table}[H]
\caption{$\Delta\varphi^*$ por ventana temporal (5 seeds por lote).
Cinco ventanas de 16 ticks cada una. Degradaci\'on = $(w_0 - w_4)/w_0$.
Las diferencias de slope no son significativas (Wilcoxon $p>0.05$).}
\label{tab:slopes}
\centering\small
\begin{tabular}{|l|c|c|c|c|}
\hline
\textbf{Ventana} & \textbf{EPR} & \textbf{EPR\_Q} & \textbf{EPR\_SHAM} & \textbf{LangGraph} \\
\hline
Ticks 0--15   & $0.490$ & $0.498$ & $0.459$ & $0.485$ \\
Ticks 16--31  & $0.484$ & $0.482$ & $0.482$ & $0.446$ \\
Ticks 32--47  & $0.455$ & $0.450$ & $0.449$ & $0.434$ \\
Ticks 48--63  & $0.452$ & $0.435$ & $0.436$ & $0.424$ \\
Ticks 64--79  & $0.441$ & $0.437$ & $0.423$ & $0.425$ \\
\hline
Degradaci\'on & $\mathbf{9.8\%}$ & $12.2\%$ & $7.8\%$ & $12.4\%$ \\
Slope $\Delta\varphi^*$ & $-0.013$ & $-0.017$ & $-0.012$ & $-0.014$ \\
\hline
\end{tabular}
\end{table}

Todos los modos muestran degradaci\'on temporal de $\Delta\varphi^*$ ---
evidencia del colapso de contexto. EPR presenta la pendiente menos
negativa ($-0.013$) con una degradaci\'on del $9.8\%$, frente a
LangGraph ($-0.014$, $12.4\%$). Sin embargo, las diferencias de slope
entre modos no son estad\'isticamente significativas (Wilcoxon
$p>0.05$), sugiriendo que el colapso de contexto es parcialmente
end\'ogeno al LLM y no es eliminado por el mecanismo de activaci\'on.
EPR\_Q degrada m\'as ($12.2\%$, slope $-0.017$): el Q-learning no
amortigua el colapso y podr\'ia amplificarlo al sobre-explotar
argumentos de alta recompensa inicial. EPR\_SHAM muestra degradaci\'on
menor ($7.8\%$) pero con d\'eficit explosivo ($3.74$): el $\delta$
aleatorio produce menos argumentos de menor calidad, reduciendo el
tama\~no del grafo y por ende el contexto acumulado susceptible de
colapsar --- no es resiliencia sino menor producci\'on. El colapso de
contexto es parcialmente end\'ogeno: el mecanismo de activaci\'on lo
modula pero no lo elimina.

\subsection{Juez LLM}

\begin{table}[H]
\caption{Puntajes del juez LLM (1--5, ciego al modo).
Evaluaci\'on offline con DeepSeek V4-Pro y GLM-5.2 en curso sobre las
20 seeds completadas.}
\label{tab:judge}
\centering\small
\begin{tabular}{|l|c|c|c|}
\hline
\textbf{Dimensi\'on} & \textbf{EPR} & \textbf{LangGraph} & $p$ (Wilcoxon) \\
\hline
Calidad argumentativa & --- & --- & --- \\
Engagement            & --- & --- & --- \\
Novelty               & --- & --- & --- \\
\hline
\end{tabular}
\end{table}

\subsection{Validaci\'on de conectividad}

\begin{table}[H]
\caption{Conectividad (5 seeds por lote).}
\label{tab:conn}
\centering\small
\begin{tabular}{|l|c|c|c|c|}
\hline
\textbf{Check} & \textbf{EPR} & \textbf{EPR\_Q} & \textbf{EPR\_SHAM} & \textbf{LangGraph} \\
\hline
Nodos aislados (post-fix)    & $0\%$ & $0\%$ & $0\%$ & $0\%$ \\
Grafo (nodos finales)        & $43.5$ & $44.9$ & $40.4$ & $60.6$ \\
Grafo (aristas finales)      & $42.4$ & $43.8$ & $39.3$ & $59.5$ \\
PRR$_G$                      & $0.971$ & $0.976$ & $0.966$ & $0.983$ \\
\hline
\end{tabular}
\end{table}

% ============================================================
\section{Discusi\'on}
\label{sec:discussion}

\paragraph{Colapso de contexto (H1).}
La degradaci\'on temporal de $\Delta\varphi^*$ mide el colapso de
contexto: todos los modos degradan, pero a distintas tasas. EPR presenta
la pendiente menos negativa ($-0.013$, degradaci\'on $9.8\%$) frente a
LangGraph ($-0.014$, $12.4\%$). Sin embargo, las diferencias de slope
entre modos no son estad\'isticamente significativas (Wilcoxon
$p>0.05$), sugiriendo que el colapso de contexto es parcialmente
end\'ogeno al LLM y no es eliminado por el mecanismo de activaci\'on.
EPR\_Q degrada m\'as ($12.2\%$, slope $-0.017$): el Q-learning no
amortigua el colapso y podr\'ia amplificarlo al sobre-explotar
argumentos de alta recompensa inicial. EPR\_SHAM muestra degradaci\'on
menor ($7.8\%$) pero con d\'eficit explosivo ($3.74$): el $\delta$
aleatorio produce menos argumentos de menor calidad, reduciendo el
tama\~no del grafo y por ende el contexto acumulado susceptible de
colapsar --- no es resiliencia sino menor producci\'on. El colapso de
contexto es parcialmente end\'ogeno: el mecanismo de activaci\'on lo
modula pero no lo elimina.

\paragraph{Distribuci\'on de turnos (H3).}
La se\~nal m\'as clara es la distribuci\'on de participaci\'on: EPR
produce max-share $0.119$ vs LangGraph $0.276$ --- una diferencia de
$2.3\times$ ($p<0.001$, Bonferroni significativo). Ning\'un agente
monopoliza el piso conversacional bajo activaci\'on homeost\'atica.
LangGraph muestra concentraci\'on sistem\'atica (max-share $0.276$).
EPR\_Q es a\'un m\'as equitativo ($0.115$), mientras que EPR\_SHAM
degrada la equidad ($0.209$, $p<0.001$). Esto sugiere que la
activaci\'on end\'ogena distribuye turnos org\'anicamente.

\paragraph{Ablaci\'on sham: el loop importa.}
EPR\_SHAM (con $\delta$ aleatorio) produce d\'eficit explosivo
($3.74$ vs $1.02$ para EPR, $p<0.001$). La forma sigmoide sola no
basta: el loop homeost\'atico ($\delta \to$ activaci\'on $\to$ calidad
$\to$ saciaci\'on) es lo que controla el d\'eficit. Esto valida que el
efecto de EPR no es artefacto de la no-linealidad del gate.
Adicionalmente, EPR\_SHAM degrada la equidad (max-share $0.209$ vs
$0.119$, $p<0.001$), confirmando que sin drive epist\'emico real la
activaci\'on se concentra.

\paragraph{Q-learning no aporta y podr\'ia perjudicar.}
EPR\_Q ($44.9$ debates, d\'eficit $0.95$) vs EPR ($43.5$ debates, d\'eficit
$1.02$): el Q-learner mejora marginalmente el control de d\'eficit pero
degrada m\'as r\'apido (slope $-0.017$ vs $-0.013$, degradaci\'on $12.2\%$
vs $9.8\%$). Las diferencias no son significativas tras Bonferroni.
La activaci\'on homeost\'atica pura es suficiente para controlar el
d\'eficit y distribuir turnos; el Q-learning a\~nade complejidad sin
beneficio claro e incluso podr\'ia amplificar el colapso de contexto al
sobre-explotar patrones iniciales exitosos.

\paragraph{LangGraph produce m\'as volumen pero menos equitativo.}
LangGraph genera $60.6$ debates vs $43.5$ de EPR, pero con mayor
concentraci\'on (max-share $0.276$ vs $0.119$, $p<0.001$) y mayor d\'eficit
m\'aximo ($2.54$ vs $1.02$, $p<0.001$). El router ex\'ogeno parece favorecer
agentes ``dominantes''. Adicionalmente, LangGraph presenta initiative ratio $=0.00$ (vs $1.00$
para EPR): ning\'un agente se activa end\'ogenamente, confirmando que la
orquestaci\'on es completamente ex\'ogena por dise\'no.

\paragraph{Amenazas a la validez.}

\textit{Interna:}
\begin{itemize}[nosep,leftmargin=*]
\item Un solo LLM (DeepSeek V4-Flash); no se puede descartar
      dependencia del proveedor.
\item RSVI est\'a integrado al locus de control, no confundido con \'el.
      RSVI verbaliza el estado regulatorio $\delta_i$ que \emph{drivea} la
      activaci\'on end\'ogena; sin el loop homeost\'atico no hay $\delta_i$
      que se actualice ni nada que verbalizar. En LangGraph, la activaci\'on
      se decide fuera del agente por un router ex\'ogeno: inyectar RSVI
      significar\'ia enviar un estado regulatorio est\'atico o ficticio que
      no drivea ninguna decisi\'on --- ruido puro, no un control v\'alido.
      La comparaci\'on es entre dos arquitecturas completas (end\'ogena con
      RSVI vs.\ ex\'ogena sin \'el), no entre ``con RSVI'' y ``sin RSVI''
      como factores independientes. Esto refleja una decisi\'on
      metodol\'ogica deliberada: evaluar el mecanismo end\'ogeno completo
      frente al paradigma ex\'ogeno establecido.
\item Los pesos del StimulusEvaluator se calibraron con seed independiente
      (grid $3{\times}3$, $\alpha\in\{3,4,5\}$, $\lambda\in\{0.02,0.03,0.05\}$);
      diferente calibraci\'on podr\'ia cambiar el comportamiento de EPR, aunque
      el perfil \emph{faction-dominant} fue seleccionado por estabilidad.
\item $g \equiv \Delta\varphi^*$ sirve simult\'aneamente como se\~nal de
      saciedad interna (AAH-3) y como m\'etrica diagn\'ostica externa.
      Esta circularidad parcial podr\'ia sesgar las m\'etricas
      estructurales automatizadas a favor de EPR, cuyo loop est\'a
      optimizado para reducir $\Delta\varphi^*$. La mitigaci\'on principal
      son los jueces LLM offline ciegos al modo, cuya evaluaci\'on es
      matem\'aticamente independiente de $g$. Sin embargo, las
      m\'etricas estructurales ($\Delta\varphi^*$, slope, AAF) deben
      interpretarse con cautela hasta confirmar concordancia con los
      jueces.
\item Resultados consolidados de 20/20 seeds (5 por lote);
      Wilcoxon signed-rank con Bonferroni ($\alpha{=}0.0056$) confirma
      diferencias significativas en distribuci\'on de turnos, d\'eficit
      y $\Delta\varphi^*$ promedio. Las pendientes temporales (slope)
      no alcanzan significancia, sugiriendo que el colapso de contexto
      es parcialmente end\'ogeno al LLM.
\end{itemize}

\textit{Estad\'istica:}
\begin{itemize}[nosep,leftmargin=*]
\item $N{=}20$ seeds pareadas; Wilcoxon signed-rank como test primario
      (no-parametrico, aprovecha el dise\~no pareado). Bonferroni con 9
      tests ($\alpha{=}0.0056$).
\item Bonferroni con 3 hip\'otesis es conservador;
      Holm-Bonferroni como sensibilidad.
\item Bootstrap 95\% CIs (10.000 resamples) para intervalos.
\item Las pendientes temporales no alcanzan significancia estad\'istica
      con $N{=}20$; se requiere mayor poder o una m\'etrica alternativa
      de colapso de contexto.
\end{itemize}

\textit{Externa:}
\begin{itemize}[nosep,leftmargin=*]
\item Estructura VSM es una de varias posibles; otras
      descomposiciones pueden producir din\'amicas distintas.
\item Un solo t\'opico de debate; efectos de t\'opico no controlados.
\item 80 heartbeats; comportamiento asint\'otico no observable.
\item Jueces LLM duales (DeepSeek V4-Pro + GLM-5.2); Cohen's $\kappa$
      para concordancia. Validaci\'on humana planificada.
\end{itemize}

\paragraph{Trabajo futuro.}
Replicaci\'on multi-LLM (GLM-5.2, Claude); reasoning effort modulado por
d\'eficit ($\delta \to$ \texttt{reasoning\_effort} en DeepSeek V4-Flash);
drive metab\'olico para regular consumo de API ($\Gamma(n{=}2)$);
dominios cooperativos; integraci\'on formal de
DeLP~\parencite{garcia2004defeasible}.

% ============================================================
\section{Conclusiones}
\label{sec:conclusion}

Presentamos un mecanismo de activaci\'on end\'ogena para sistemas
multi-agente de LLM basado en homeostasis epist\'emica, formalizado
mediante los Axiomas de Autonom\'ia Homeost\'atica (AAH). La condici\'on
EPR ($\Gamma(n{=}1)$, sin Q-learning) a\'isla el aporte del mecanismo
regulatorio b\'asico: sigmoide, d\'eficit y sensor estructural. El
dise\~no experimental con 20 semillas pareadas (4 lotes de 5), 4 condiciones (EPR,
EPR\_Q, EPR\_SHAM, LangGraph), hiperpar\'ametros calibrados con seed
independiente, y jueces LLM duales ciegos al modo, permite aislar el
efecto del locus de control. Los resultados (20/20 seeds, Wilcoxon
signed-rank con Bonferroni) muestran que EPR distribuye turnos
$2.3\times$ m\'as equitativamente que LangGraph (max-share $0.119$ vs
$0.276$, $p<0.001$), controla el d\'eficit ($1.02$ vs $2.54$,
$p<0.001$) y produce mayor $\Delta\varphi^*$ promedio ($0.463$ vs
$0.442$, $p<0.001$). La ablaci\'on sham confirma que el loop
homeost\'atico --- no s\'olo la forma del gate --- es lo que controla el
d\'eficit ($3.74$ vs $1.02$, $p<0.001$). El Q-learning no aporta
beneficio claro y podr\'ia amplificar el colapso. Las pendientes
temporales no alcanzan significancia estad\'istica, sugiriendo que el
colapso de contexto es parcialmente end\'ogeno al LLM y no es eliminado
por el mecanismo de activaci\'on. La evaluaci\'on con jueces LLM
offline ciegos al modo est\'a en curso.

\printbibliography

\end{document}
